\begin{lemma}[Independent coefficients from redundant coordinates]
\label{lem:selmer-padded-relations}
Let $T$ be a finite-dimensional vector space over a field $F$, let
$U\subseteq T^g$, and set $W=\{t:e_jt\in U\text{ for all }j\}$.
If $\operatorname{codim}U\leq g$ and $\dim W+g^2\leq D$, there are
$D$ forms $X_a$ spanning $T^*$ and $g^2$ independent vectors
$c_{ij}\in F^D$ such that
\[
             \sum_j\sum_a c_{ij,a}X_a(u_j)=0
             \qquad(u\in U,\ 1\leq i\leq g).
\]
\end{lemma}

\begin{proof}
Choose $g$ spanning equations $\Phi_i$ for $U$, allowing zero equations,
and write $\Phi_i(u)=\sum_j\phi_{ij}(u_j)$. The common kernel of the
$g^2$ forms $\phi_{ij}$ is $W$. A complement to their span in $T^*$
therefore has dimension $\dim W$. List all the $\phi_{ij}$, then
a basis of that complement, and pad with zero forms to obtain $D$
forms $X_a$. Choose the $c_{ij}$ to be the distinct coordinate
vectors selecting the listed $\phi_{ij}$. Their independence holds
even if those forms are dependent or zero. For $g=0$, only a spanning
family of $T^*$ is needed.
\end{proof}

In the deformation presentation, $U$ consists of the tuples of even
tangent classes that admit the prescribed second-order lift. The
second cohomological condition gives $\operatorname{codim}U\leq g$.
The first bounds the independent classes furnished by $W$ and the
$g^2$ classes $\varepsilon_j\kappa_i$, giving $\dim W+g^2\leq D$.
Thus the independence in Lemma~\ref{lem:selmer-local-incidence}
concerns the presenting power-series ring; no additional independence
of the images of its coordinates is required.

\begin{lemma}[Separated blocks and a uniform local bound]
\label{lem:selmer-block-resultant}
Let $A$ be a discrete valuation ring with uniformizer $t$ and finite
residue field $k$. Let $P=\operatorname{diag}(B,C)$, with block sizes
$a,b$, and suppose
\[
 \delta=\operatorname{Res}(\det(XI_a-B),\det(XI_b-C))\ne0.
\]
An odd matrix, meaning one with zero diagonal blocks, that commutes
with $P$ over $A$ is zero. For every $m\geq0$,
\[
 \#\{Z\text{ odd over }A/t^m:PZ=ZP\}\leq |k|^{2v_t(\delta)}.
\]
\end{lemma}

\begin{proof}
On the two off-diagonal blocks the commutator is $U\mapsto BU-UC$
and $V\mapsto CV-VB$. Triangularizing the blocks over an algebraic
closure of the fraction field shows that the determinant of their
direct sum is $(-1)^{ab}\delta^2$. In particular its kernel over $A$
is zero. If its elementary divisors have valuations $s_r$, its kernel
modulo $t^m$ has order $|k|^{\sum_r\min(m,s_r)}$, at most
$|k|^{\sum_rs_r}=|k|^{2v_t(\delta)}$. Empty blocks give a zero odd
module and $\delta=1$.
\end{proof}

For the auxiliary Frobenius matrices this gives the single constant
$2\sum_jv_t(\delta_j)$ for the added odd local terms, independent of
the truncation and of sufficiently large auxiliary level. In rank
three with block sizes $1,2$, it is enough to use
$\delta_j=\det(\alpha_jI_2-C_j)$. Eigenvalues inside one block may
coincide; eigenvalue differences across the blocks need not be units.
Other uses of the auxiliary primes retain their stated hypotheses.

The same calculation proves descent at the limiting group. Commuting
Frobenius and tame elements $\varphi_j,\tau_j$, with $\tau_j$ in the
kernel of the map to the base, force an odd cocycle's value
$z(\tau_j)$ to commute with $P_j$, hence to vanish. Its zero set
in that kernel is closed and normal under the index-two subgroup.
The conjugate-dual identity makes it stable under complex conjugation
as well. Since the $\tau_j$ normally generate the kernel topologically,
the cocycle descends and is spanned, after inverting $t$, by the original
odd Selmer classes. Kernel generation and descent of the prescribed
local conditions are retained hypotheses of this application.

\begin{lemma}[A bounded obstruction]
\label{lem:selmer-finite-obstruction}
Let $H$ be a finite-length module over a discrete valuation ring $A$,
killed by $t^m$, and let $g,C,m$ be nonnegative integers. If
\[
          \operatorname{length}_A(H/t^eH)\leq ge+C
          \qquad(0\leq e\leq m),
\]
then any $g+1$ elements $\xi_i\in H$ satisfy
$t^C\sum_i c_i\xi_i=0$ for coefficients $c_i\in A$ with at least
one $c_i$ a unit.
\end{lemma}

\begin{proof}
If $C\geq m$, the assertion is immediate. Otherwise write
$H=\bigoplus_r A/(t^{\lambda_r})$. At $e=C+1$ the inequality shows
that at most $g$ of the $\lambda_r$ exceed $C$. Thus $t^CH$ is
generated by at most $g$ elements. Lift the $g+1$ elements $t^C\xi_i$
to $A^g$. Their matrix has a nonzero kernel vector over the fraction
field. Clearing denominators and dividing by the common smallest power
of $t$ gives a kernel vector over $A$ with a unit coordinate.
\end{proof}

The remaining arithmetic estimates must provide one additive constant
for all truncations, allowing the required auxiliary level to depend
on the truncation. The local bound above supplies its auxiliary odd
term. The global and ordinary local counts, and the construction of
the patched automorphic modules, remain necessary.
